%=================================================
% Corpo do relatório
%=================================================

\section{Introdução}

O Aprendizado por Reforço trata de agentes que aprendem a decidir por meio da interação direta com o ambiente, ajustando seu comportamento a partir das recompensas que recebem \cite{sutton2018reinforcement}. Diferentemente do planejamento em um Processo de Decisão de Markov (MDP) conhecido, o agente não dispõe das probabilidades de transição nem da função de recompensa, de modo que precisa estimá-las, ainda que implicitamente, a partir das tuplas $(s, a, r, s')$ que observa ao atuar \cite{russell2016artificial}. Nesse contexto, o algoritmo Q-Learning se destaca por aprender diretamente os valores ótimos das ações, sem exigir um modelo do ambiente.

Neste trabalho, o problema escolhido é o do robô coletor de lixo, inspirado no robô Robby descrito em \cite{mitchell2009complexity}. O agente, chamado Marvin, percorre uma grade bidimensional de $50 \times 50$ células com paredes internas e lixo distribuído aleatoriamente, enxergando somente a célula em que está e as quatro células adjacentes. Como consequência, Marvin precisa aprender uma política de coleta sem saber onde está e sem conhecer o mapa.

Além da implementação do simulador e do Q-Learning tabular, o trabalho investiga experimentalmente o compromisso entre exploração e explotação, que é o ponto central da diferença entre planejar e aprender. Para isso, comparam-se diferentes valores fixos de $\epsilon$ e diferentes agendas de decaimento, mantendo os demais hiperparâmetros constantes.

O restante do relatório está organizado da seguinte forma. A Seção~\ref{sec:modelagem} descreve o ambiente e formaliza o problema como um MDP. Em seguida, a Seção~\ref{sec:metodologia} apresenta o algoritmo e a escolha dos hiperparâmetros, enquanto a Seção~\ref{sec:resultados} reúne as curvas de aprendizado e os experimentos de exploração. Por fim, a Seção~\ref{sec:politica} analisa a política aprendida, e as Seções~\ref{sec:discussao} e~\ref{sec:conclusao} discutem as limitações da modelagem e concluem o trabalho.

%=================================================
\section{Ambiente e modelagem do problema}
\label{sec:modelagem}

\subsection{Descrição do ambiente}

O ambiente é uma grade de $50 \times 50$ células, na qual cada célula pode conter lixo ($+1$), espaço vazio ($0$) ou parede ($-1$). As coordenadas $(x, y)$ seguem a convenção de indexação de matrizes, com $x$ indicando a coluna, $y$ indicando a linha e a origem $(0,0)$ no canto superior esquerdo. Marvin inicia todo episódio em $s_{\text{start}} = (0,0)$, que é garantidamente livre de parede. Além disso, qualquer tentativa de sair da grade é tratada como colisão com parede, o que mantém Marvin na mesma célula e gera a penalidade correspondente.

Cada episódio tem duração fixa de $H = 1000$ passos. Ao término, o ambiente é reinicializado, com o lixo redistribuído e Marvin de volta a $s_{\text{start}}$. A geração da grade segue densidades de \pendente{densidade de paredes} de paredes e \pendente{densidade de lixo} de lixo, sorteadas uniformemente.

\pendente{Figura com um recorte do ambiente gerado, salva em resultados/}

\subsection{Espaço de estados}

Marvin não observa a própria posição, de modo que a representação por $(x, y)$ não é adequada. Além de não ser observável, ela tampouco é Markoviana, pois a melhor ação em uma célula depende de onde ainda há lixo. Por isso, adota-se como estado a percepção local do agente,
\begin{equation}
    s = \left(c_{\text{atual}},\, c_N,\, c_S,\, c_L,\, c_O\right),
    \label{eq:estado}
\end{equation}
em que $c_{\text{atual}} \in \{\text{vazio}, \text{lixo}\}$ e $c_N, c_S, c_L, c_O \in \{\text{vazio}, \text{lixo}, \text{parede}\}$. Como a célula atual nunca é parede, o número de estados é
\begin{equation}
    |S| = 2 \cdot 3^4 = 162,
\end{equation}
e a Q-table possui $162 \times 7 = 1134$ entradas. Assumiu-se ainda que o limite da grade é percebido pelos sensores como parede, o que é coerente com o tratamento de colisão adotado pelo ambiente \pendente{confirmar premissa com os professores}.

Embora compacta e observável, essa representação sofre de ambiguidade perceptual, uma vez que células distintas da grade podem gerar a mesma percepção. Consequentemente, o processo percebido pelo agente não é estritamente Markoviano, aspecto retomado na Seção~\ref{sec:discussao}.

\subsection{Ações}

O conjunto de ações é
\begin{equation}
    A(s) = A = \{\text{norte}, \text{sul}, \text{leste}, \text{oeste}, \text{aleatória}, \text{parado}, \text{coletar}\},
\end{equation}
com $|A| = 7$. Esse conjunto não depende de $s$ porque qualquer ação pode ser tentada em qualquer estado. Por exemplo, andar contra uma parede ou coletar em uma célula sem lixo não são proibidos, mas apenas penalizados. Dessa maneira, cabe ao agente aprender, por meio da recompensa, quais ações fazem sentido em cada percepção.

\subsection{Função de recompensa}

A recompensa depende do par $(s, a)$ e do resultado da ação, e não apenas do estado de chegada. Em notação de MDP, tem-se
\begin{equation}
    R(s, a, s') =
    \begin{cases}
        +10 & \text{se } a = \text{coletar e } c_{\text{atual}} = \text{lixo}, \\
        -1  & \text{se } a = \text{coletar e } c_{\text{atual}} \neq \text{lixo}, \\
        -5  & \text{se o movimento é inválido (parede ou fora da grade)}, \\
        0   & \text{nos demais casos.}
    \end{cases}
    \label{eq:recompensa}
\end{equation}
Para as quatro direções, o movimento é inválido quando a célula vizinha correspondente é parede, informação que está contida na percepção. Já para a ação aleatória, a penalidade depende da direção sorteada. Ademais, o fato de $s' = s$ não implica colisão, pois uma célula vizinha diferente pode produzir a mesma percepção. Por esse motivo, a recompensa é calculada pelo ambiente a partir da grade real.

\subsection{Transições}

Das sete ações, somente a ação aleatória é estocástica. Denotando por $d$ a direção sorteada uniformemente em $\{\text{norte}, \text{sul}, \text{leste}, \text{oeste}\}$, a transição é dada por
\begin{equation}
    T(s, \text{aleatória}, s') = \frac{1}{4} \sum_{d} T(s, d, s').
\end{equation}
Para as demais ações, o efeito sobre a célula é determinístico, ou seja, $T(s, a, s') = 1$ para a célula resultante e $0$ para as demais. Cabe ressaltar, contudo, que a percepção seguinte $s'$ também depende de células que o agente não observa, de modo que, do ponto de vista do agente, as transições não são determinísticas na prática.

\subsection{Fator de desconto e horizonte}

Adotou-se $\gamma = 0{,}9$ \pendente{validar por teste de sensibilidade}. Como o lixo é um objetivo local, um horizonte efetivo de aproximadamente $1/(1-\gamma) = 10$ passos é suficiente. Em contrapartida, valores de $\gamma$ próximos de 1 propagam ruído por muito tempo, enquanto valores muito baixos tornariam o agente míope.

O problema é episódico com horizonte fixo $H = 1000$ e não possui estado terminal. Por isso, a função $\text{IsEnd}(s)$ não faz sentido neste caso, já que o episódio termina por limite de tempo, e não por Marvin atingir um estado especial. Como consequência, no último passo de cada episódio a Eq.~\eqref{eq:atualizacao} mantém o termo de bootstrap, pois o corte é artificial e o valor futuro de $s'$ continua existindo, embora não seja observado.

%=================================================
\section{Metodologia}
\label{sec:metodologia}

\subsection{Q-Learning tabular}

A cada transição observada $(s, a, r, s')$, a estimativa do valor ótimo da ação é atualizada por
\begin{equation}
    \hat{Q}^{\star}(s, a) \leftarrow (1 - \alpha)\, \hat{Q}^{\star}(s, a) + \alpha \left[ r + \gamma \max_{a' \in A(s')} \hat{Q}^{\star}(s', a') \right],
    \label{eq:atualizacao}
\end{equation}
em que $\alpha$ é a taxa de aprendizado \cite{sutton2018reinforcement}. \pendente{Complementar com a forma de erro de diferença temporal e a explicação do alvo e da predição}

\subsection{Política de atuação $\epsilon$-greedy}

\pendente{Descrever a Eq. (4) e a implementação}

\subsection{Hiperparâmetros}

\pendente{Tabela com $\alpha$, $\gamma$, $Q_0$, $K$, $H$, sementes e justificativas}

%=================================================
\section{Resultados}
\label{sec:resultados}

\subsection{Curvas de aprendizado}

\pendente{Recompensa por episódio com média móvel de 20; lixos coletados, coletas malsucedidas e movimentos inválidos; fração da Q-table visitada}

\subsection{Exploração versus explotação}

\pendente{Tabela da varredura de $\epsilon$ fixo e discussão dos extremos $\epsilon = 0$ e $\epsilon = 1$}

\pendente{Comparação das agendas constante, exponencial e harmônica no mesmo gráfico}

%=================================================
\section{Política aprendida}
\label{sec:politica}

\pendente{Mapeamento estado para ação: estados com lixo na célula atual e estados cercados por paredes}

%=================================================
\section{Discussão e limitações}
\label{sec:discussao}

\pendente{Ambiguidade perceptual, ausência de memória, dependência das densidades escolhidas, número de execuções agregadas}

%=================================================
\section{Conclusão}
\label{sec:conclusao}

\pendente{Síntese dos resultados e trabalhos futuros}